%! Author = <autor tłumaczenia>
%! Kompilacja: lualatex "<plik>.tex" (2-3 razy, dla spisu treści)
\documentclass[10pt,twoside,a4paper]{article}
\usepackage{<projekt>-style}   % kopia assets/gry-style.sty
\graphicspath{{./png/<wersja>/}{./png/}}

\renewcommand{\seriatytul}{<Nazwa serii / gry jak w paginie oryginału>}
\renewcommand{\stopkatekst}{<Tytuł instrukcji> v<wersja> · tłumaczenie nieoficjalne}

\title{<Tytuł> -- <podtytuł PL>}
\author{<autor tłumaczenia>}

\begin{document}

% --- Okładka (ilustracja z oryginału, bez logo wydawcy) ------------------------
\begin{titlepage}
  \thispagestyle{empty}
  \begin{tikzpicture}[remember picture,overlay]
    \node[anchor=south,inner sep=0pt] at (current page.south)
      {\includegraphics[width=\paperwidth,height=0.78\paperheight,keepaspectratio=false]{cover}};
    \fill[top color=wgziel!75!black,bottom color=wgziel!70!white]
      (current page.north west) rectangle ([yshift=-0.22\paperheight]current page.north east);
    \node[anchor=north,align=center,text=white,font=\naglowekfont\bfseries]
      at ([yshift=-14mm]current page.north)
      {\fontsize{34}{38}\selectfont <Tytuł linia 1>\\[2pt]
       \fontsize{34}{38}\selectfont <Tytuł linia 2>\\[8pt]
       \fontsize{24}{28}\selectfont <Podtytuł PL>\\[6pt]
       \fontsize{24}{28}\selectfont Wersja <wersja>};
    \node[anchor=south east,align=right,fill=white,fill opacity=.82,text opacity=1,
      inner sep=6pt,font=\naglowekfont\small]
      at ([xshift=-8mm,yshift=8mm]current page.south east)
      {Nieoficjalne tłumaczenie na język polski\\<autor tłumaczenia>};
  \end{tikzpicture}
\end{titlepage}

% --- Wstęp / przedmowa (poza multicols; plik sam zarządza kolumnami) ----------
\import{./tex/}{00_powitanie.tex}

% --- Spis treści (zwarty, ma się zmieścić na 1 stronie) -------------------------
\clearpage
\begin{multicols}{2}
  \fontsize{8.6}{10}\selectfont
  \setlength{\cftbeforesecskip}{2pt}
  \tableofcontents
\end{multicols}
\clearpage

% --- Zasady (wszystkie rozdziały w jednym multicols) -----------------------------
\begin{multicols}{2}
  \import{./tex/}{01_wstep.tex}
  % \import{./tex/}{02_....tex}
  % ...
\end{multicols}

% --- Dodatki -----------------------------------------------------------------
\clearpage
\import{./tex/}{15_tabele.tex}       % tabele z ostatniej strony oryginału (poza multicols)
\clearpage
\import{./tex/}{14_slowniczek.tex}   % generowany: scripts/gen_glossary.py

\end{document}
